\section{Confidence intervals from sampling behaviour}

Suppose a statistic \(T\) estimates a parameter \(\theta\), and assume that
\[
\frac{T-\theta}{s_T}\approx N(0,1),
\]
where \(s_T\) is known and does not depend on \(\theta\). Then
\[
P\left(T-1.96s_T\leq\theta\leq T+1.96s_T\right)\approx0.95.
\]
The direct algebra is valid because \(s_T\) is fixed.

\begin{definitionbox}
A \textbf{confidence interval} is produced by a procedure designed to contain the
fixed parameter with a stated long-run frequency under repeated sampling.
\end{definitionbox}

After the sample is observed, the interval is numerical. The confidence level
describes the long-run coverage of the procedure, not a probability distribution
for the fixed parameter.

\subsection*{Mean}

For independent observations with mean \(\mu\) and variance \(\sigma^2\), a
large-sample interval with known \(\sigma\) is
\[
\overline X\pm1.96\frac{\sigma}{\sqrt n}.
\]
When \(\sigma\) is unknown, the large-sample plug-in version is
\[
\overline X\pm1.96\frac{S}{\sqrt n},
\qquad
S^2=\frac1{n-1}\sum_{i=1}^n(X_i-\overline X)^2.
\]

\subsection*{Proportion}

For a sample proportion, the theoretical standard error
\[
\sqrt{\frac{p(1-p)}{n}}
\]
depends on the unknown \(p\). Replacing it by
\[
\sqrt{\frac{\widehat p(1-\widehat p)}{n}}
\]
gives the approximate interval
\[
\widehat p\pm1.96\sqrt{\frac{\widehat p(1-\widehat p)}{n}}.
\]
This is a substitution-based approximation, not the literal algebraic inversion
of the theoretical expression. It can be inaccurate for small samples or values
near the endpoints.

\subsection*{Interpretation}

Larger samples usually reduce standard errors at the rate \(1/\sqrt n\), but they
do not remove bias, dependence, systematic measurement error, or an inappropriate
model.

\begin{summarybox}
Direct inversion is straightforward when the standard error is known and does
not depend on the unknown parameter. Estimated standard errors require an
additional approximation that should be stated explicitly.
\end{summarybox}
